\makeabstract
{
Understanding Mania: The Role of Emotion Beliefs and Emotion Regulation in Mania Symptoms​
}
{
Tiffany Sun (1), Mabel Shanahan (1), Elizabeth T. Kneeland (1)
}
{
\textsuperscript{1}Department of Psychology, Amherst College, Amherst, MA, USA
}
{
This study explores the influence of emotion beliefs and emotion regulation strategies on mania symptoms. Previous research indicates that viewing emotions as fixed, unique, or long-lasting is linked to greater psychopathology, and maladaptive cognitive strategies specifically increase risk for mania. By exploring the role these factors have in mania symptoms, the study aims to improve our understanding of mania and offer insights that could inform effective approaches to managing mania symptoms.
}
{
In this study, 186 participants (66.1\% women; mean age = 28; 61.9\% White/Caucasian) completed several questionnaires through Qualtrics that assessed their baseline mania symptoms, general emotion beliefs, and emotion regulation strategies. The latter two were then reassessed through one-month follow-up surveys. Since the study examines mania specifically, only data from 41 participants who reported elevated mania symptoms indicative of a highly probable manic or hypomanic condition were included in the analysis. 
}
{
Baseline correlation analysis revealed that mania symptom severity is positively associated with maladaptive regulation strategies such as brooding rumination, catastrophizing, and beliefs that emotions are unique and long-lasting. Interestingly, mania symptom severity is not significantly associated with emotion malleability beliefs. Further analysis into brooding rumination revealed that brooding thoughts of feeling wronged, self-criticism, or perceiving one’s experiences as unique are linked to greater mania severity. At the one-month follow-up, longitudinal analysis indicated that baseline mania severity was positively correlated with stronger brooding thoughts of self-blame and perceiving experiences as unique.
}
{
This study concludes that certain negative emotion beliefs and maladaptive regulation strategies are linked to greater mania symptom severity. Contrary to expectations, emotion malleability beliefs are not significantly associated with mania symptoms. These findings improve our understanding of mania and offer insights toward developing more effective strategies in managing mania symptoms. Future research could expand these findings with a larger, clinical sample and a longitudinal design that measures mania severity at multiple time points.
}

\newpage
